% Class-compatibility preamble retained for reader excerpts of
%   Hiraku Nakajima, "Quiver Varieties and Branching", SIGMA 5 (2009) 003,
%   DOI 10.3842/SIGMA.2009.003, CC BY 4.0.
%
% The source pins the publisher class sigma, which is not part of the delivery
% baseline, so this file supplies the amsart baseline and then copies the author's
% own declaration block mechanically and unchanged (pinned-original-source.tex
% lines 43-239): the aenume environment, the Theorem/Corollary/Lemma/Proposition and
% definition-style newtheorem declarations, the equation numbering, the author's
% reference macros (\lemref and the like) and the author's mathematical macros.
% Everything from pinned source line 240 onwards is left out: those lines define
% \rf, which uses the publisher class internal dimension \y@boxdim, and the two
% colour rules that serve that one macro; nothing in the delivered excerpts uses
% \rf, \graysquare or \halfgray.
% No mathematical macro, statement, symbol or assumption is introduced, changed or
% removed by this file.  The author's package mathdots (pinned source line 20) is not
% reproduced because it is not installed in the delivery TeX installation and no
% delivered excerpt uses its commands; the author's remaining packages are kept.
\documentclass[10pt,a4paper]{amsart}
\usepackage{amssymb}
\usepackage{amsbsy,amscd,mathtools}
\usepackage{mathrsfs}
\usepackage{nicefrac}
\usepackage[all]{xy}
\usepackage{color}
\usepackage{hyperref}
\newenvironment{aenume}{%
  \begin{enumerate}%
  \renewcommand{\theenumi}{\alph{enumi}}%
  \renewcommand{\labelenumi}{\textup{(\theenumi)}}%
  }{\end{enumerate}}
%

\usepackage{verbatim}
\usepackage{version}
\usepackage{color}
\newenvironment{NB}{
\color{red}{\bf NB}. \footnotesize
}{}
% For a final version, NB disappears.
\excludeversion{NB}
% Older version
%\renewenvironment{NB}{%
%  \comment
%  }{\endcomment}
% The following is used inside NB
\newenvironment{NB2}{
\color{blue}{\bf NB}. \footnotesize
}{}
%       Theorem environments
%% \theoremstyle{plain} %% This is the default
\newtheorem{Theorem}[equation]{Theorem}
\newtheorem{Corollary}[equation]{Corollary}
\newtheorem{Lemma}[equation]{Lemma}
\newtheorem{Proposition}[equation]{Proposition}
%\newtheorem{Axiom}{Axiom}

\theoremstyle{definition}
\newtheorem{Definition}[equation]{Definition}
\newtheorem{Example}[equation]{Example}
\newtheorem{CounterExample}[equation]{Counter-Example}
\newtheorem{Notation}[equation]{Notation}
\newtheorem{Convention}[equation]{Convention}
\newtheorem{Conjecture}[equation]{Conjecture}

%\theoremstyle{remark}
\newtheorem{Remark}[equation]{Remark}
\newtheorem{Remarks}[equation]{Remarks}
%\newtheorem*{Claim}{Claim}
%\renewcommand{\theClaim}{} % to make the claim environment unnumbered
\newtheorem{Question}[equation]{Question}
%\renewcommand{\theQuestion}{}  % to make the notation environment
                               % unnumbered
\newtheorem{Exercise}[equation]{Exercise}
%\renewcommand{\theExercise}{}

\numberwithin{equation}{section}
%\numberwithin{Theorem}{section}
%\numberwithin{figure}{section}

\newcommand{\thmref}[1]{Theorem~\ref{#1}}
\newcommand{\secref}[1]{Section~\ref{#1}}
\newcommand{\lemref}[1]{Lemma~\ref{#1}}
\newcommand{\propref}[1]{Proposition~\ref{#1}}
\newcommand{\corref}[1]{Corollary~\ref{#1}}
\newcommand{\subsecref}[1]{Section~\ref{#1}}
\newcommand{\chapref}[1]{Chapter~\ref{#1}}
\newcommand{\defref}[1]{Definition~\ref{#1}}
\newcommand{\remref}[1]{Remark~\ref{#1}}
\newcommand{\remsref}[1]{Remarks~\ref{#1}}
%
%       Math definitions
%\newcommand{\lsp}[2]{{}^{#1}{#2}}
\newcommand{\lsp}[2]{{\mskip-.3mu}{}^{#1}\mskip-1mu{#2}}
\newcommand{\defeq}{\overset{\operatorname{\scriptstyle def}}{=}}
\newcommand{\interior}{\rule{6pt}{.4pt}\rule{.4pt}{6pt}} % interior
                                                         % product
\newcommand{\C}{{\mathbb C}}
\newcommand{\Z}{{\mathbb Z}}
\newcommand{\Q}{{\mathbb Q}}
\newcommand{\R}{{\mathbb R}}
% Projective spaces
\newcommand{\proj}{{\mathbb P}}
%\newcommand{\CP}{\operatorname{\C P}}
\newcommand{\CP}{\proj}
\newcommand{\RP}{\operatorname{\R P}}
% Lie Groups
\newcommand{\SL}{\operatorname{\rm SL}}
\newcommand{\SU}{\operatorname{\rm SU}}
\newcommand{\GL}{\operatorname{GL}}
\newcommand{\PGL}{\operatorname{PGL}}
\newcommand{\U}{\operatorname{\rm U}}
\newcommand{\SO}{\operatorname{\rm SO}}
\newcommand{\grpSp}{\operatorname{\rm Sp}}
% Lie algebras
\newcommand{\algsl}{\operatorname{\mathfrak{sl}}} % because \sl="slant"
\newcommand{\su}{\operatorname{\mathfrak{su}}}
\newcommand{\gl}{\operatorname{\mathfrak{gl}}}
\newcommand{\algu}{\operatorname{\mathfrak{u}}}
\newcommand{\so}{\operatorname{\mathfrak{so}}}
\newcommand{\algsp}{\operatorname{\mathfrak{sp}}}
\newcommand{\g}{{\mathfrak g}}
\newcommand{\h}{{\mathfrak h}}
% Useful symbols
\newcommand{\Spec}{\operatorname{Spec}\nolimits}
\newcommand{\Proj}{\operatorname{Proj}\nolimits}
\newcommand{\End}{\operatorname{End}}
\newcommand{\Hom}{\operatorname{Hom}}
\newcommand{\Ext}{\operatorname{Ext}}
\newcommand{\Ker}{\operatorname{Ker}}
\newcommand{\Coker}{\operatorname{Coker}}
\newcommand{\Ima}{\operatorname{Im}}
\newcommand{\coker}{\operatorname{coker}}
\newcommand{\Endskew}{\operatorname{Endskew}}
\newcommand{\im}{\mathop{\text{\rm im}}\nolimits}
%\newcommand{\dim}{\mathop{\text{\rm dim}}\nolimits}
\newcommand{\rank}{\operatorname{rank}}
\newcommand{\ind}{\mathop{\text{\rm ind}}\nolimits}
\newcommand{\tr}{\operatorname{tr}}
\newcommand{\Tr}{\operatorname{Tr}}
\newcommand{\DB}{\overline{\partial}}
\newcommand{\pd}[2]{\frac{\partial#1}{\partial#2}}
\newcommand{\deriv}[1]{\frac{d}{d#1}}
\newcommand{\id}{\operatorname{id}}
\newcommand{\ve}{\varepsilon}
%
% notations used in this paper
%
% \newcommand{\vin}[1]{\operatorname{in}(#1)} % incoming vertex
% \newcommand{\vout}[1]{\operatorname{out}(#1)} % outgoing vertex
\newcommand{\vin}[1]{\operatorname{i}(#1)} % incoming vertex
\newcommand{\vout}[1]{\operatorname{o}(#1)} % outgoing vertex
\newcommand{\bM}{{\mathbf M}} % the vector space
\newcommand{\bN}{{\mathbf N}} % the vector space
\newcommand{\M}{{\mathfrak M}} % the moduli space
\newcommand{\Mreg}{\M^{\operatorname{s}}}
\newcommand{\La}{{\mathfrak L}} % the Lagrangian variety
\newcommand{\Fi}{{\mathfrak F}} % the fixed point variety
\newcommand{\dslash}{/\!\!/} % for algebro-geometric quotient (double slash)
\newcommand{\bv}{{\mathbf v}} % the vector v
\newcommand{\bw}{{\mathbf w}} % the vector w
\newcommand{\bC}{{\mathbf C}} % the Cartan matrix
\newcommand{\bA}{{\mathbf A}} % the adjacency matrix
\newcommand{\bI}{{\mathbf I}} % the identity matrix
\newcommand{\bn}{{\mathbf n}} % the kernel of the extended Cartan matrix
\newcommand{\bu}{{\mathbf u}}
\newcommand{\be}{{\mathbf e}}
\newcommand{\bq}{{\mathbf q}}
\newcommand{\bs}{\mathbf s}
\newcommand{\bc}{\mathbf c}
\newcommand{\codim}{\operatorname{codim}} % codimension
\newcommand{\gr}{\operatorname{gr}} % grade
\newcommand{\topdeg}{{\operatorname{top}}} % top degree
\newcommand{\Uq}{{\mathbf U}_q(\mathfrak g)} % the QUE algebra
\newcommand{\Un}{{\mathbf U}_q}
\newcommand{\Um}{\Tilde{\mathbf U}_q(\mathfrak g)} % the modified
                                                 % enveloping algebra
\newcommand{\Ua}{{\mathbf U}_q(\widehat{\mathfrak g})} % the quantum affine
                                                 % algebra
\newcommand{\Uae}{{\mathbf U}_{\varepsilon}(\widehat{\mathfrak g})} %
                                                 % the quantum affine
                                                 % algebra
\newcommand{\Ul}{{\mathbf U}_q({\mathbf L}{\mathfrak g})} % the quantum
                                                      % loop algebra
\newcommand{\Uma}{\Tilde{\mathbf U}^\bw_q(\widehat{\mathfrak g})}
%                                                 the quantum affine algebra
\newcommand{\Ui}{{\mathbf U}^{\Z}_q(\mathfrak g)}
\newcommand{\Ue}{{\mathbf U}_{\varepsilon}(\mathfrak g)}
\newcommand{\Ueone}{{\mathbf U}_{\varepsilon = 1}(\mathfrak g)}
\newcommand{\Uli}{{\mathbf U}^{\Z}_q({\mathbf L}\mathfrak g)}
\newcommand{\Ule}{{\mathbf U}_{\varepsilon}({\mathbf L}\mathfrak g)}
\newcommand{\bU}{\mathbf U} % the enveloping algebra
\newcommand{\Uqm}{{\widetilde{\mathbf U}_q}} % its modified form
\newcommand{\Pa}{{\mathfrak P}} % Hecke correspondence
\newcommand{\HomE}{\operatorname{E}}
\newcommand{\HomL}{\operatorname{L}}
\newcommand{\Zreg}{Z^{\operatorname{reg}\nolimits}}
\newcommand{\hiw}{{L}} % the highest weight module
%      \interval is used to provide better spacing after a [ that
%      is used as a closing delimiter.
\newcommand{\interval}[1]{\mathinner{#1}}
\newcommand{\shfO}{\mathcal O}
\newcommand{\Mp}{\widetilde\M}
\newcommand{\Mpp}{\widehat\M}
\newcommand{\Wedge}{{\textstyle \bigwedge}}
%
\newcommand{\Zw}{Z(\bw)}
\newcommand{\Mw}{\M(\bw)}
\newcommand{\Law}{\La(\bw)}
%\newcommand{\Zw}{Z(\bullet,\bullet;\bw)}
%\newcommand{\Mw}{\M(\bullet,\bw)}
%\newcommand{\Law}{\La(\bullet,\bw)}
\newcommand{\ngpiso}{\overset{.}{=}}
\newcommand{\idl}{\mathfrak I}
\newcommand{\Ktop}{K_{\operatorname{top}}}
\newcommand{\Kztop}{K_{0,\operatorname{top}}}
\newcommand{\Kotop}{K_{1,\operatorname{top}}}

\newcommand{\qH}{\mathbb H}
\newcommand{\D}{\mathcal D}
\newcommand{\A}{\mathcal A}
\newcommand{\incarrow}{\xrightarrow{\mathrm{inclusion}}}
\newcommand{\proarrow}{\xrightarrow{\mathrm{projection}}}
\usepackage{cleveref}
\crefname{originallemma}{original lemma}{original lemmas}
\Crefname{originallemma}{Original lemma}{Original lemmas}
\makeatletter
\newlabel{lem:betasurj}{{2.18}{source}{}{sourcepointer-1}{}}
\newlabel{lem:betasurj@cref}{{[originallemma][0][]2.18}{[1][0][]source}{}{}{}}
\makeatother
\crefname{originaltheorem}{original theorem}{original theorems}
\Crefname{originaltheorem}{Original theorem}{Original theorems}
\makeatletter
\newlabel{thm:CB}{{2.11}{source}{}{sourcepointer-2}{}}
\newlabel{thm:CB@cref}{{[originaltheorem][0][]2.11}{[1][0][]source}{}{}{}}
\makeatother
\crefname{originallemma}{original lemma}{original lemmas}
\Crefname{originallemma}{Original lemma}{Original lemmas}
\makeatletter
\newlabel{lem:face}{{2.8}{source}{}{sourcepointer-3}{}}
\newlabel{lem:face@cref}{{[originallemma][0][]2.8}{[1][0][]source}{}{}{}}
\makeatother
\title{The associated graded of a zeta-stable module is a direct sum of simple modules when the module is stable for a generic parameter}
\author{Hiraku Nakajima}
\date{doi:10.3842/SIGMA.2009.003; CC BY 4.0}
\begin{document}
\maketitle
\noindent\textit{Editorial note.} All delivered mathematical text is the author's own, copied from the pinned publisher source sources/algebra-fl01-quiver-varieties-and-branching/pinned-original-source.tex; no statement or proof step was rewritten, reordered or removed.\par\medskip
\noindent\textit{Editorial note.} Preamble provenance: as for the other cards of this source, the retained preamble supplies the amsart baseline in place of the publisher class sigma and copies the author's declaration block of pinned source lines 43-239 unchanged.\par\medskip
\noindent\textit{Editorial note.} Bibliography provenance: the counted proof cites the author's reference Na-proc. That citation is resolved with sources/algebra-fl01-quiver-varieties-and-branching/bibliography.tex, an excerpt of the author's own thebibliography environment (pinned source lines 4098-4257) copied character for character; no entry was rewritten and no citation key was changed.\par\medskip
\noindent\textit{Editorial note.} Pairing provenance: the author's proof carries the explicit header Proof of Lemma lem:bulletstable (source line 1963), so the statement at 1921-1936 and the proof at 1963-2026 are paired by the author's own words; the remark at 2028-2032 that follows the proof is a separate remark and is not counted.\par\medskip
\noindent\textit{Editorial note.} External references are navigation pointers only; every deliberately kept step of the author's argument stays inside the counted spans. Their printed numbers are the author's own numbering under the author's newtheorem declarations, computed mechanically by the lane's scanner.\par\medskip
\noindent\textit{Allowed standard prerequisites (editorial).}\begin{itemize}
\item The zeta-bullet-stability of a module (B,a,b), the parameter zeta-bullet and the notion of a generic parameter in the sense of the paper. Defined earlier in the paper (the author's Definition at source lines 746-787 and the subsections on stability, source lines 726-1019). Fixes the hypothesis of the statement and the parameter with respect to which the module is assumed stable.
\item The face structure of the set of stability conditions and the decomposition of a polystable module into its stable components. Used inside the counted proof, where the author argues through the faces of the stability parameter and the associated graded object. Supplies the device by which the associated graded of a polystable module is described and why the summands are simple.
\end{itemize}\medskip
\ifdefined\hypertarget\hypertarget{sourcepointer-1}{}\fi
\noindent\textit{Source reference.} Original lemma 2.18: sources/algebra-fl01-quiver-varieties-and-branching/pinned-original-source.tex, source lines 1710--1710. Navigation pointer only: the counted proof appeals to this lemma for the surjectivity of the maps beta\_i; the lemma's own statement and proof are not reproduced here and the property it supplies is declared as a named standard prerequisite.\par\medskip
\ifdefined\hypertarget\hypertarget{sourcepointer-2}{}\fi
\noindent\textit{Source reference.} Original theorem 2.11: sources/algebra-fl01-quiver-varieties-and-branching/pinned-original-source.tex, source lines 1219--1219. Navigation pointer only: the counted proof uses the author's earlier non-emptiness criterion in one step to recognise which simple modules can occur; the criterion itself is not reproduced.\par\medskip
\ifdefined\hypertarget\hypertarget{sourcepointer-3}{}\fi
\noindent\textit{Source reference.} Original lemma 2.8: sources/algebra-fl01-quiver-varieties-and-branching/pinned-original-source.tex, source lines 999--999. Navigation pointer only: the counted proof refers to the author's face-structure lemma in one step of the face argument; the lemma itself is not reproduced here.\par\medskip
\begin{Lemma}\label{lem:bulletstable}
  A $\zeta^\bullet$-stable point $B\in\mu^{-1}(0)\subset\bM(\bv,0)$
  satisfying the normalization condition $\zeta^\bullet\cdot\bv = 0$
  is of one of the following three forms\textup:
  \begin{aenume}\itemsep=0pt
  \item a point in $X_{\zeta^\circ}\setminus \bigcup_{i\in I_0^0}
    C_i$,
  \item $\bv$ is a coordinate vector $\be_i$ for an $i\in
    I_0^0$, or
  \item $\bv$ is $\delta - \alpha_h$, where $\alpha_h$ is the highest
    root vector of a connected component of the sub-Dynkin diagram
    $I_0^0$.
  \end{aenume}
  In cases \textup{(b)}, \textup{(c)} the corresponding
  $\zeta^\bullet$-stable point is unique up to isomorphisms.
\end{Lemma}

\begin{proof}[Proof of \lemref{lem:bulletstable}]
  Consider the criterion in \thmref{thm:CB}. Since our graph is of
  af\/f\/ine type, we have $p(x) = 1$ if $x$ is an imaginary root (i.e.,
  $x = m\delta$ with $m\ge 1$) and $p(x) = 0$ if $x$ is a real root.
  Therefore if $\bv$ is imaginary root, the condition is equivalent to
  $\bv = \delta$. This is the case (a). If $\bv$ is a real root, then
  the condition is that $\bv$ cannot have a nontrivial decomposition
  $\bv = \sum_t \beta^{(t)}$ with $\zeta^\bullet\cdot \beta^{(t)} = 0$. Then
  it is clear that we have either (b) or (c).


Next show that a $\zeta^\bullet$-stable point is unique in case (b),
(c). It is clear in the case (b). Consider the case (c). We apply an
argument used by Mukai for the case of rigid sheaves on a $K3$
surface:
%
Let $B$, $B'$ be $\zeta^\bullet$-stable points with the underlying
$I$-graded vector space $V$ of dimension vector $\delta-\alpha_h$. We
consider the complex
\begin{gather*}
  \HomL(V,V) \xrightarrow{\alpha} \HomE(V,V) \xrightarrow{\beta}
  \HomL(V,V),
\\
  \alpha(\xi) = B\xi - \xi B', \qquad
  \beta(C) = \varepsilon (B C - C B').
\end{gather*}
Since the alternating sum of dimensions is $2$, we cannot have both
$\alpha$ is injective and $\beta$ is surjective. If $\alpha$ is not
injective, we have a nonzero homomorphism $\xi$ from the module $B$ to
$B'$. Then considering submodules $\Ker\xi\subset B$ and
$\Ima\xi\subset B'$, the $\zeta^\bullet$-stability of $B$ and $B'$
imply $\Ker\xi = 0$ and $\Ima\xi = B'$. Therefore $B$ and $B'$ are
isomorphic. If $\beta$ is not surjective, we consider
$\beta^*\in \HomL(V^*,V^*) \cong \HomL(V,V)$ instead, we get the same
assertion.

Finally consider the case $\bv = \delta$.
%
By \lemref{lem:face}(3b) and the subsequent remark $B$ is
$\zeta^\circ$-stable. So $B$ corresponds to a point $[B]$ in
$X_{\zeta^\circ}$. Therefore it only remains to determine which point
in $X_{\zeta^\circ}$ is $\zeta^\bullet$-stable.

If $B$ is not $\zeta^\bullet$-stable, we consider a
Jordan--H\"older f\/iltration to f\/ind a proper quotient which is
$\zeta^\bullet$-stable. By the above discussion, the quotient must
have the dimension vector $\be_i$ with $i\in I_0^0$ or
$\delta-\alpha_h$. But in the latter case, $B$ contains a submodule
with the dimension vector $\be_i$ with $i\in I_0^0$. This is
impossible for the $\zeta^\circ$-stability, as $\zeta^\circ \cdot
\be_i > 0$.
%
Thus $B$ has the quotient with the dimension vector $\be_i$ with
$i\in I_0^0$. In particular, $\beta$ above is not surjective, and
hence $B$ is contained in the exceptional component $C_i$ by
\cite[Proof of 5.10]{Na-proc}.
%
On the other hand, if $B$ is in $C_i$, $\beta$ is not surjective,
and it has a quotient with $\dim = \be_i$. It is not
$\zeta^\bullet$-stable.
%
Therefore $\zeta^\bullet$-stable points are precisely in
$X_{\zeta^\circ}\setminus \bigcup_{i\in I_0^0} C_i$.
\end{proof}

% Bibliography excerpt retained for reader excerpts of
%   Hiraku Nakajima, "Quiver Varieties and Branching", SIGMA 5 (2009) 003,
%   DOI 10.3842/SIGMA.2009.003, CC BY 4.0.
%
% Every entry below is copied character for character from the author's own
% thebibliography environment of the pinned publisher source
% (pinned-original-source.tex lines 4098-4257).  No entry was rewritten,
% reordered, merged, split or invented, and no citation key was changed;
% only the entries that the delivered excerpts cite are kept.
\begin{thebibliography}{99}

\bibitem{CB}
Crawley-Boevey W., Geometry of the moment map for representations of
  quivers, {\it Compositio Math.} \textbf{126} (2001),   257--293.
 % \MR{MR1834739 (2002g:16021)}


\bibitem{Kac}
 Kac V.G., Inf\/inite-dimensional Lie algebras, 3rd ed., Cambridge
  University Press, Cambridge, 1990. %\MR{MR1104219 (92k:17038)}


\bibitem{KS}
Kashiwara M., Saito Y., Geometric construction of crystal bases, {\it Duke
  Math. J.} \textbf{89} (1997),  9--36,  \mbox{\href{http://arxiv.org/abs/q-alg/9606009}{q-alg/9606009}}. %\MR{MR1458969 (99e:17025)}


\bibitem{Na-quiver}
Nakajima H., Instantons on {ALE} spaces, quiver varieties, and
  Kac--Moody algebras, {\it Duke Math. J.} \textbf{76} (1994), 365--416.
  %\MR{MR1302318 (95i:53051)}


\bibitem{Na-proc}
Nakajima H., Varieties associated with quivers, in Representation Theory of
  Algebras and Related Topics (Mexico City, 1994), {\it CMS Conf. Proc.}, Vol.~19,
  Amer. Math. Soc., Providence, RI, 1996, 139--157.
  %\MR{MR1388562   (97m:16022)}


\bibitem{Na-alg}
Nakajima H., Quiver varieties and Kac--Moody algebras, {\it Duke Math. J.}
  \textbf{91} (1998), 515--560. %\MR{MR1604167 (99b:17033)}


\bibitem{Lecture}
Nakajima H., Lectures on Hilbert schemes of points on surfaces, {\it University
  Lecture Series}, Vol.~18, American Mathematical Society, Providence, RI, 1999.
  %\MR{MR1711344 (2001b:14007)}


\bibitem{Na-qaff}
Nakajima H., Quiver varieties and f\/inite-dimensional representations of
  quantum af\/f\/ine algebras, {\it J. Amer. Math. Soc.} \textbf{14} (2001),
  145--238, \href{http://arxiv.org/abs/math.QA/9912158}{math.QA/9912158}. %\MR{MR1808477 (2002i:17023)}


\bibitem{Na-tensor}
Nakajima H., Quiver varieties and tensor products, {\it Invent. Math.}
  \textbf{146} (2001), 399--449, \href{http://arxiv.org/abs/math.QA/0103008}{math.QA/0103008}. %\MR{MR1865400 (2003e:17023)}


\bibitem{Na-reflect}
Nakajima H.,  Ref\/lection functors for quiver varieties and {W}eyl group
  actions, {\it Math. Ann.} \textbf{327} (2003), 671--721.
  %\MR{MR2023313   (2004k:16036)}


\bibitem{Na-ADHM}
Nakajima H.,  Sheaves on $ALE$ spaces and quiver varieties, {\it Mosc. Math. J.}
  \textbf{7} (2007),  699--722.


\bibitem{Saito}
Saito Y., Crystal bases and quiver varieties, {\it Math. Ann.} \textbf{324}
  (2002),  675--688, \href{http://arxiv.org/abs/math.QA/0111232}{math.QA/0111232}. %\MR{MR1942245 (2004a:17023)}



\end{thebibliography}
\end{document}
